\subsection{FRM EXAM 2007---QUESTION 50 (Bonds): duration-based VaR}

\textbf{Enunciado}

\emph{A portfolio consists of two zero coupon bonds, each with a current value of \$10. 
The first bond has a modified duration of one year and the second has a modified duration 
of nine years. The yield curve is flat, and all yields are 5\%. Assume all moves of the 
yield curve are parallel shifts. Given that the daily volatility of the yield is 1\%, 
which of the following is the best estimate of the portfolio’s daily value at risk (VAR) 
at the 95\% confidence level?}
\begin{enumerate}[label=\alph*.]
\item \hl{USD 1.65}
\item USD 2.33
\item USD 1.16
\item USD 0.82
\end{enumerate}

\subsubsection*{Solución}

\paragraph{1. Definiciones} Sean: 

\begin{itemize}
    \item \(P\): precio actual del bono.
    \item \(\Delta P\): cambio en el precio del bono durante el día.
    \item \(y\): rendimiento al vencimiento (\emph{yield}).
    \item \(\Delta y\): cambio diario en el rendimiento.
    \item \(D\): duración modificada.
\end{itemize}

\vspace{0.5cm}

Entonces:
\[
\frac{\Delta P}{P}
\]
representa el cambio porcentual en el precio del bono.\\

\paragraph{2. Duración modificada}

La duración modificada se define como
\[
D \;=\; -\frac{1}{P}\frac{dP}{dy}.
\]
Que se obtiene a partir de:
\[
P(y)=\frac{F}{(1+y)^T}.
\]

Derivando respecto de \(y\). Esta definición dice que \(D\) mide qué tan sensible es el precio del bono a cambios
pequeños en la yield, se puede interpretar que dividir entre \(P\) normaliza por el precio, eso convierte una sensibilidad en moneda a una sensibilidad porcentual. 

Haciendo álgrebra de la NASA,
\[
\frac{dP}{P} \;=\; -D\,dy.
\]

Para movimientos pequeños de tasas, reemplazamos diferenciales por cambios finitos
pequeños y obtenemos la aproximación lineal:
\[
\frac{\Delta P}{P}\approx -D\,\Delta y.
\]

Multiplicando ambos lados por \(P\),
\[
\Delta P \approx -D\,P\,\Delta y.
\]

\paragraph{3. Interpretación de la fórmula}

La expresión
\[
\Delta P \approx -D\,P\,\Delta y
\]
significa lo siguiente:

\begin{itemize}
    \item si la tasa sube (\(\Delta y>0\)), el precio del bono baja (\(\Delta P<0\));
    \item si la tasa baja (\(\Delta y<0\)), el precio del bono sube (\(\Delta P>0\));
    \item mientras mayor sea \(D\), más sensible es el bono a cambios en tasa;
    \item mientras mayor sea \(P\), mayor será el cambio en dólares.
\end{itemize}

Así, \(D\,P\) es la exposición en una moneda \(P\) del bono al factor de riesgo
\(\Delta y\).

\paragraph{4. Aplicación a cada bono}

El portafolio tiene dos bonos cupón cero, cada uno con precio actual
\[
P_1=P_2=10.
\]

Sus duraciones modificadas son
\[
D_1=1,
\qquad
D_2=9.
\]

Como el problema dice que todos los movimientos de la curva son
\textbf{desplazamientos paralelos}, ambos bonos enfrentan el mismo cambio
\(\Delta y\).

Entonces, para el bono 1:
\[
\Delta P_1 \approx -D_1P_1\Delta y = -(1)(10)\Delta y = -10\Delta y.
\]

Para el bono 2:
\[
\Delta P_2 \approx -D_2P_2\Delta y = -(9)(10)\Delta y = -90\Delta y.
\]

\paragraph{5. Cambio del valor del portafolio}

El cambio total del portafolio es la suma de ambos cambios:
\[
\Delta P_{\text{port}}=\Delta P_1+\Delta P_2.
\]

Sustituyendo:
\[
\Delta P_{\text{port}}
\approx -10\Delta y-90\Delta y
= -100\Delta y.
\]

Esto muestra que el portafolio tiene una exposición total de \$100 por unidad de cambio en el rendimiento.

\paragraph{6. ¿Cómo entra la volatilidad del yield?}

El problema dice que la volatilidad diaria del rendimiento es
\[
\sigma_y=1\%=0.01.
\]

Como
\[
\Delta P_{\text{port}}\approx -100\Delta y,
\]
entonces la desviación estándar del cambio diario en el valor del portafolio es
\[
\sigma_{\Delta P}
=100\,\sigma_y
=100(0.01)
=1.
\]

Es decir, el cambio diario del valor del portafolio tiene una volatilidad
aproximada de \$1.

\paragraph{7. Cálculo del VaR al 95\%}

Bajo la aproximación paramétrica, el VaR diario al 95\% se estima como
\[
\mathrm{VaR}_{0.95}\approx z_{0.95}\,\sigma_{\Delta P}.
\]

Para una normal estándar,
\[
z_{0.95}\approx 1.645.
\]

Por tanto,
\[
\mathrm{VaR}_{0.95}\approx 1.645(1)=1.645.
\]

Redondeando:
\[
{\mathrm{VaR}_{0.95}\approx \$1.65.}
\]